%!TEX root = TheoretischePhysikC.tex

\section{Vorlesung: Entropie}
Wie wird der Begriff ,,gleich wahrscheinlich verteilte Mikrozustände`` in der Quantenmechanik definiert? Um diese Frage zu beantworten müssen wir das Konzept von \emph{Entropie} einführen. Diese Größe wird als Maß für die \emph{Informationsinhalt} eines Zustandes betrachtet. Oder als ein Maß an \emph{fehlender Information}. Je mehr Informationsgehalt ein Zustand besitzt, desto weniger wissen wir, welcher Mikrozustand auftritt.

Drei Lesarten laufen nebeneinander her: fehlende Information, Unordnung, Informationsgehalt -- dieselbe Zahl, verschieden gelesen. Nötig wird sie, weil im Hilbertraum das klassische Abzählen der Mikrozustände versagt: wegen der Superpositionen gibt es unendlich viele.

\begin{intuition}
	Ein Speicher fasse zehn Bit. Kennen wir seinen Inhalt nicht, so fehlen uns zehn Bit; stehen alle Bits auf Null, fehlt uns keine Information mehr -- und es ist auch kein Speicherplatz mehr frei.
\end{intuition}

Sei $\mathcal{H} \equiv \mathbb{C}^d$ ein Hilbertraum und $\rho$ ein Dichteoperator, der das statistische Ensemble $\{(p_j,\rho_j)\,|\, j=1,2, \ldots, N\}$ darstellt:
\begin{equation}
  \rho \equiv \sum_{j=1}^N p_j \rho_j.
\end{equation}

Die von Neumann-Entropie von $\rho$ wird wie folgt definiert
\begin{equation}
  \eqbox{S(\rho) = -k\tr(\rho\log(\rho)) = -k\langle \log(\rho)\rangle,}    
\end{equation} 
wobei $k=1,38064852 \times 10^{-23} \, \text{m}^2 \,\text{kg} \,\text{s}^{-2} \,\text{K}^{-1}$ die Boltzmann-Konstante ist und $\log(\rho)$ wird mit Hilfe vom Funktionalkalkül berechnet. Der Logarithmus ist hier, wie überall in der Physik, der natürliche; und $\log(\rho)$ muss selbst ein Operator sein, sonst gäbe es das Produkt $\rho\log(\rho)$ gar nicht.

Weil jeder Dichteoperator $\rho$ hermitesch ist, können wir $\rho$ bezüglich einer Eigenwert-Eigenvektor-Zerlegung darstellen:
\begin{equation}
  \rho = \sum_{j=1}^d r_j |r_j\rangle\langle r_j|,
\end{equation}
wobei $r_j\ge 0$ und $\langle r_j|r_k\rangle = \delta_{jk}$. Der Logarithmus von $\rho$ wird als
\begin{equation}
  \log(\rho) = \sum_{j=1}^d \log(r_j) |r_j\rangle\langle r_j|
\end{equation}
berechnet. Ist ein Eigenwert $r_j = 0$, so ist $\log(r_j) = -\infty$; darauf kommen wir im ersten Beispiel zurück.

\begin{merksatz}[Vorsicht mit dem Logarithmus einer Matrix]
	$\log(\rho)$ ist \emph{nicht} der eintragsweise Logarithmus der Einträge von $\rho$. Erst diagonalisieren, dann den Logarithmus auf die Eigenwerte anwenden. Eintragsweise rechnen darf man nur, wenn $\rho$ bereits diagonal ist.
\end{merksatz}

Ein Beispiel: für
\begin{equation}
	\rho = \begin{pmatrix} \frac{1}{2} & \frac{1}{4}\\[2pt] \frac{1}{4} & \frac{1}{2}\end{pmatrix}
	\quad \text{ist} \quad
	\log(\rho) = \frac{1}{2}\begin{pmatrix} \log\frac{3}{16} & \log 3\\[2pt] \log 3 & \log\frac{3}{16}\end{pmatrix}
	\neq \begin{pmatrix} \log\frac{1}{2} & \log\frac{1}{4}\\[2pt] \log\frac{1}{4} & \log\frac{1}{2}\end{pmatrix},
\end{equation}
denn die Eigenwerte sind $3/4$ und $1/4$.

Daher wird die von Neumann-Entropie von $\rho$ wie folgt berechnet:
\begin{equation}
  S(\rho) = -k\sum_{j=1}^d \log(r_j) \tr(|r_j\rangle\langle r_j|\rho) = -k\sum_{j=1}^d r_j\log(r_j).
\end{equation}
Das ist der Befreiungsschlag: die Entropie hängt \emph{nur} von den Eigenwerten ab, die Eigenvektoren tauchen nicht mehr auf. Eine unitäre Drehung ändert das Spektrum nicht, also auch nicht die Entropie:
\begin{equation}
	\eqbox{S(\rho) = S(U\rho U^{\dagger}) \quad \text{für alle unitären } U.}
\end{equation}

\begin{übung}
	Beweisen Sie die unitäre Invarianz $S(\rho) = S(U\rho U^{\dagger})$.
\end{übung}

\begin{beispiel}
  Ein reiner Zustand $|\psi\rangle$ entspricht einem statischen Ensemble, in dem $|\psi\rangle$ mit Wahrscheinlichkeit $1$ auftritt. Der dazugehörige Dichteoperator ist durch
  \begin{equation}
    \rho = |\psi\rangle\langle \psi|
  \end{equation}
  gegeben. Dieser Operator befindet sich bereits in diagonaler Form wobei $r_1 = 1$ (und $r_2 = r_3 = \cdots = r_d =0$). Daher lautet die von Neumann-Entropie
  \begin{equation}
    S(\rho) = -k\sum_{j=1}^d r_j\log(r_j) = -k r_1\log(r_1) = 0.
  \end{equation}
  Hierbei nimmt die Funktion $f(x) = -x\log(x)$ den Wert $0$ an der Stelle $x=0$ an. Genauer ist $0\log(0)$ als Grenzwert $\lim_{x\to 0} x\log(x)$ zu lesen; die Kurve $-x\log(x)$ läuft bei $x=0$ sichtbar auf Null zu. In diesem Beispiel beträgt die von Neumann-Entropie Null. Die Interpretation ist hier, dass das statistische Ensemble keine Informationen enthält. Man weiß immer, in welchem Zustand sich das System befindet. Es fehlt keine Information über den Zustand des Systems. Ein reiner Zustand ist eben kein Ensemble.
\end{beispiel}

\begin{beispiel}
  Angenommen, $\mathcal{H}$ ist der Hilbertraum $\mathbb{C}^d$. Nehmen wir weiter an, dass $\{|j\rangle \}_{j=1}^d$ eine Orthonormalbasis für $\mathcal{H}$ ist. Der vollständig gemischte Zustand $\rho$ ist gegeben durch
  \begin{equation}
    \rho \equiv \sum_{j=1}^d \frac{1}{d} |j\rangle\langle j|.
  \end{equation}
  Dieser Zustand entspricht einem statischen Ensemble, in dem $|j\rangle$ mit gleich verteilter Wahrscheinlichkeit $1/d$ auftritt. Dieser Operator befindet sich auch bereits in diagonaler Form wobei $r_1 = 1/d = r_2 = \cdots = r_d$. Die von Neumann-Entropie wird wie folgt berechnet:
  \begin{equation}
    S(\rho) = -k\sum_{j=1}^d r_j\log(r_j) = -k \times d \times  (1/d) \times \log(1/d) = k \log(d).
  \end{equation}
  Die von Neumann-Entropie erreicht in diesem Fall den maximal möglichen Wert. Die Interpretation ist hier, dass das statistische Ensemble die maximale Information enthält. Man weiß nicht, in welchem Zustand sich das System befindet. Es fehlen alle Information über den Zustand des Systems.
\end{beispiel}

% fig_vl19_entropie.tex — Die von-Neumann-Entropie als Summe der Terme
% -r log r: der einzelne Term, und die Entropie eines Qubits.
%
% WHY THIS FIGURE.  The lecture states three things it never shows: that
% -x log x runs into zero at x = 0 (the text says the curve does so
% "sichtbar", so the curve had better be there — it is the one picture
% Osborne draws at the board, v23 [64:14]), that S = 0 for a pure state
% and S = k log d at the other extreme, and that only the spectrum
% enters.  What prose cannot hand over is the SHAPE between those two
% extremes: that S rises with infinite slope out of the pure state — a
% one-percent admixture already buys eight percent of the maximum — and
% then flattens off, so that near the maximally mixed state the entropy
% hardly notices how the weight is distributed.  That shape, and the
% fact that the whole reachable range of a qubit is the strip below
% k log 2, is the content of the plate.
%
% GEOMETRY FROM THE FORMULA.  Nothing is drawn by eye.  Both curves are
% the lecture's own expression S(rho) = -k sum_j r_j log(r_j) evaluated
% on a spectrum, in units of k:
%
%   (a)  eta(r) = -r log r          one eigenvalue's contribution,
%   (b)  S/k    = eta(p) + eta(1-p) a qubit with spectrum {p, 1-p}.
%
% Every marked value is exact, not eyeballed:
%   * eta has its maximum at r = 1/e, of height 1/e = 0.36788 — the
%     y tick and the grey guide are placed there;
%   * eta(1/2) = (1/2) log 2 = 0.34657, and twice that is the peak of
%     panel (b): the red dot in (a) is literally half the red dot in
%     (b).  That is why both are red — same physical situation, two
%     eigenvalues of equal weight;
%   * S/k at p = 1/2 is log 2 = 0.69315 = log d for d = 2, the ceiling
%     of the lecture's second example;
%   * S/k at p = 3/4 is -(3/4)log(3/4) - (1/4)log(1/4) = 0.56234, the
%     value of the worked example in the text, whose two matrices
%     (1/4)diag(3,1) and (1/4)[[2,1],[1,2]] share the spectrum {3/4,1/4}
%     and therefore land on the SAME dot.  Drawing them arriving at one
%     point is the unitary invariance S(rho) = S(U rho U^dagger) shown
%     rather than restated;
%   * the dotted rule between p = 1/4 and p = 3/4 is the same statement
%     one level down: a spectrum is a set, so p and 1-p are one state
%     and must have one entropy.
%
% DESIGN DECISIONS
%   * colour ramps with the entropy value, and carries nothing else:
%     open circle (spacecadet on white) where the term vanishes —
%     r = 0 never occurs, r = 1 occurs for certain, neither hides any
%     information; birrenorange for the intermediate example; birrenred
%     for the maximum.  The curves themselves are the house data colour
%     cgblue in both panels because they are the same function;
%   * the band above log 2 is hatched, not filled: it is unreachable for
%     d = 2, and hatching survives greyscale printing.  The grey rule
%     bounding it is a reference line, so it is neutral, like the guide
%     to the maximum of eta;
%   * no y labels: the panel titles already name the plotted quantity,
%     and the width saved goes into the annotations, which sit on their
%     objects rather than in a legend;
%   * the words under the p axis ("rein", "maximal gemischt") name the
%     states, not the numbers — the reading of the figure is physical;
%   * annotation positions were set against the 12pt amsart build, not
%     against scripts/tikz-preview.sh: the preview is a 10pt standalone,
%     so every \scriptsize label is noticeably smaller there and
%     clearances that look ample in the PNG can close up on the page.
%     The eta(1/2) label therefore sits above the curve, centred on
%     r = 1/2 and tied to its dot by a hairline leader (to the right of
%     the dot the label grew into panel (b); a rule carried down to the
%     axis instead crossed the words of the limit annotation), and the
%     annotation block of (b) is centred on p = 1/2, the axis of symmetry
%     of the curve, which is where the free space is widest;
%   * grey does two jobs and they are kept apart: dashed rules are
%     levels (1/e, log 2), solid hairlines are leaders from a label to
%     the thing it names.

\begin{center}
\begin{tikzpicture}

% ─────────────────────────────────────────────────────────────
%  Exact constants (see header)
% ─────────────────────────────────────────────────────────────
\def\einve{0.3678794}       % 1/e             — maximum of eta
\def\halflogtwo{0.3465736}  % (1/2)log 2      — eta(1/2)
\def\logtwo{0.6931472}      % log 2           — S/k of the mixed qubit
\def\Sdreiviertel{0.5623351}% S({3/4,1/4})/k  — the worked example

% ─────────────────────────────────────────────────────────────
%  (a)  The single term eta(r) = -r log r
% ─────────────────────────────────────────────────────────────
\begin{axis}[
    name=term,
    tpcplot,
    width=5.8cm, height=5.4cm,
    grid=none, clip=false,
    xmin=-0.05, xmax=1.12, ymin=-0.035, ymax=0.50,
    xtick={0,0.5,1}, xticklabels={$0$,$\tfrac12$,$1$},
    ytick={0,\einve}, yticklabels={$0$,$1/e$},
    xlabel={$r$},
    title={(a) \tpcifenglish{one term}{ein Term}: $\eta(r)=-r\log r$},
    domain=0.0002:1, samples=300,
]
  % One reference line only: the height of the maximum.  The equal-weight
  % point gets a leader down from its label instead of a rule through the
  % whole panel — a full-height rule at r = 1/2 runs straight through the
  % words of the limit annotation below, and dashed grey is reserved here
  % for level lines, solid hairlines for leaders.
  \draw[levellinegrey] (axis cs:-0.05,\einve) -- (axis cs:\einve,\einve);

  \addplot[wavefn, cgblue] {-x*ln(x)};

  % Both zeros of eta: nothing is missing there
  \node[circle, draw=spacecadet, fill=white, line width=0.7pt,
        inner sep=1.1pt] at (axis cs:0,0) {};
  \node[circle, draw=spacecadet, fill=white, line width=0.7pt,
        inner sep=1.1pt] at (axis cs:1,0) {};

  % Equal weights: half of the qubit maximum in (b)
  \node[circle, fill=birrenred, inner sep=1.55pt]
    at (axis cs:0.5,\halflogtwo) {};
  \node[lbl, font=\scriptsize, anchor=south]
    at (axis cs:0.5,0.405) {$\eta(\tfrac12)=\tfrac12\log 2$};
  \draw[birrengray, line width=0.4pt]
    (axis cs:0.5,0.399) -- (axis cs:0.5,0.357);

  % The limit that defines 0 log 0
  \node[lbl, font=\scriptsize, anchor=west]
    at (axis cs:0.085,0.055) {$\lim_{r\to0} r\log r = 0$};
  \draw[birrengray, line width=0.4pt]
    (axis cs:0.078,0.047) -- (axis cs:0.022,0.013);

  % The state names hang a fixed 6.2 mm below the axis line, here and in
  % (b).  In axis units the same drop would land at two different heights
  % on the page: the panels cover 0.535 and 0.945 units in the same
  % 5.4 cm, so one unit is nearly twice as long in (a) as in (b).
  \node[lbl, text=birrengray, font=\scriptsize, anchor=north, yshift=-6.2mm]
    at (axis cs:0,-0.035) {\tpcifenglish{never}{nie}};
  \node[lbl, text=birrengray, font=\scriptsize, anchor=north, yshift=-6.2mm]
    at (axis cs:1,-0.035) {\tpcifenglish{certain}{sicher}};
\end{axis}

% ─────────────────────────────────────────────────────────────
%  (b)  A qubit: S/k = eta(p) + eta(1-p)
% ─────────────────────────────────────────────────────────────
\begin{axis}[
    name=qubit,
    at={(term.south east)}, anchor=south west, xshift=12mm,
    tpcplot,
    width=6.8cm, height=5.4cm,
    grid=none, clip=false,
    xmin=-0.04, xmax=1.07, ymin=-0.045, ymax=0.90,
    xtick={0,0.25,0.5,0.75,1},
    xticklabels={$0$,$\tfrac14$,$\tfrac12$,$\tfrac34$,$1$},
    ytick={0,\logtwo}, yticklabels={$0$,$\log 2$},
    xlabel={$p$},
    title={(b) \tpcifenglish{qubit}{Qubit}: $S/k=\eta(p)+\eta(1-p)$},
    domain=0.0004:0.9996, samples=320,
]
  % Unreachable for d = 2: hatched, bounded by a neutral reference rule
  \fill[forbidden] (axis cs:-0.04,\logtwo) rectangle (axis cs:1.07,0.90);
  \draw[levellinegrey] (axis cs:-0.04,\logtwo) -- (axis cs:1.07,\logtwo);
  % the backing is `tpcplot`'s legend style, but opaque: at the legend's
  % 0.88 the hatching still shows through the words
  \node[lbl, text=birrengray, font=\scriptsize, inner sep=2.5pt,
        fill=white, rounded corners=1.5pt]
    at (axis cs:0.5,0.795)
    {\tpcifenglish{unreachable}{unerreichbar}: $S_{\max}=k\log d$};

  \addplot[wavefn, cgblue] {-x*ln(x)-(1-x)*ln(1-x)};

  % One spectrum, two places on the p axis
  \draw[classicallim] (axis cs:0.25,\Sdreiviertel)
                   -- (axis cs:0.75,\Sdreiviertel);
  \node[circle, fill=birrenorange, inner sep=1.55pt]
    at (axis cs:0.25,\Sdreiviertel) {};
  \node[circle, fill=birrenorange, inner sep=1.55pt]
    at (axis cs:0.75,\Sdreiviertel) {};

  % Maximum and the two pure states
  \node[circle, fill=birrenred, inner sep=1.55pt]
    at (axis cs:0.5,\logtwo) {};
  \node[circle, draw=spacecadet, fill=white, line width=0.7pt,
        inner sep=1.1pt] at (axis cs:0,0) {};
  \node[circle, draw=spacecadet, fill=white, line width=0.7pt,
        inner sep=1.1pt] at (axis cs:1,0) {};

  % The two matrices of the worked example arrive at the same dot
  \node[lbl, font=\footnotesize, anchor=north] (mats) at (axis cs:0.5,0.44)
    {$\tfrac14\!\left(\begin{smallmatrix}3&0\\0&1\end{smallmatrix}\right)$
     \tpcifenglish{and}{und}
     $\tfrac14\!\left(\begin{smallmatrix}2&1\\1&2\end{smallmatrix}\right)$};
  \draw[-{Stealth[length=4pt]}, birrengray, line width=0.5pt,
        shorten >=2.5pt] (mats.north east) -- (axis cs:0.75,\Sdreiviertel);
  \node[lbl, text=birrengray, font=\scriptsize, anchor=north, align=center]
    at (axis cs:0.5,0.315)
    {\tpcifenglish{same spectrum,}{gleiches Spektrum,}\\
     \tpcifenglish{same entropy}{gleiche Entropie}};

  \node[lbl, text=birrengray, font=\scriptsize, anchor=north, yshift=-6.2mm]
    at (axis cs:0,-0.045) {\tpcifenglish{pure}{rein}};
  \node[lbl, text=birrengray, font=\scriptsize, anchor=north, yshift=-6.2mm]
    at (axis cs:1,-0.045) {\tpcifenglish{pure}{rein}};
  \node[lbl, text=birrengray, font=\scriptsize, anchor=north, yshift=-6.2mm]
    at (axis cs:0.5,-0.045)
    {\tpcifenglish{max.\ mixed}{maximal gemischt}};
\end{axis}

\end{tikzpicture}
\end{center}


Die Abbildung zeigt links den einzelnen Term $\eta(r)=-r\log r$: er verschwindet
an beiden Enden -- bei $r=0$ tritt der Zustand nie auf, bei $r=1$ mit Sicherheit.
Rechts steht die Summe zweier solcher Terme, die Entropie eines Qubits. Achten Sie
auf die steile Flanke: schon ein Prozent Beimischung zum reinen Zustand kostet acht
Prozent des Maximums $k\log 2$; die beiden Matrizen mit demselben Spektrum landen
dagegen auf einem einzigen Punkt.

Wegen des Vorfaktors $k$ ist die Entropie eines einzelnen Qubits winzig; groß wird sie erst bei sehr vielen Teilchen.

\begin{übung}\label{ex:entropie-additiv}
	Es sei $\rho = \rho_1\otimes\rho_2\otimes\cdots\otimes\rho_N$ der Zustand von $N$ Teilchen. Zeigen Sie die Additivität
	\begin{equation}
		S(\rho) = \sum_{j=1}^N S(\rho_j).
	\end{equation}
	Für gleiche Faktoren ist das $N$-fache der Entropie eines einzelnen Teilchens -- die Entropie ist \emph{extensiv}.
\end{übung}